\documentclass[border=5mm]{standalone}
% ==============================================================================
% SETUP.TEX - CONFIGURACIÓN GLOBAL DEL PROYECTO
% ==============================================================================
% Este archivo centraliza todos los paquetes y estilos.
% Debe incluirse en el preámbulo del libro principal y de los archivos standalone.
% \PassOptionsToPackage{gray}{xcolor}
% --- 1. CODIFICACIÓN E IDIOMA ---
% fix-cm habilita las fuentes Computer Modern en tamaños arbitrarios (por debajo
% de 5 pt, entre otros). Debe cargarse antes que fontenc.
\usepackage{fix-cm}
\usepackage[utf8]{inputenc}
\usepackage[T1]{fontenc}
\usepackage[spanish, es-tabla]{babel}  
\usepackage{csquotes} % Recomendado para babel y citas

\usepackage{import}
% mode=image: Usa el PDF pre-compilado, nunca intenta recompilar.
% Debes compilar cada figura manualmente antes de compilar el libro.
\usepackage[mode=image, subpreambles=false]{standalone}

% --- 2. MATEMÁTICAS Y CIENCIAS ---
\usepackage{amsmath, amssymb, amsfonts, mathtools}
\usepackage{enumitem}

% LaTeX trae \DeclareMathSizes{5}{5}{5}{5}, así que a 5 pt (\tiny) los subíndices
% salen del mismo tamaño que la base: en las etiquetas de figuras el M de
% $\omega_M$ queda desproporcionado. Se restablece la proporción habitual.
\DeclareMathSizes{5}{5}{3.7}{3}

% --- 3. DISEÑO Y GEOMETRÍA --- 
% Unificamos los márgenes para todo el documento
\usepackage[left=2.5cm,right=2.5cm,top=3cm,bottom=3cm]{geometry}
\makeatletter
\ifstandalone % Evita que geometry dañe el recorte de las figuras bajándolas 3cm
  \@ifclasswith{standalone}{crop=false}{
    % Es un capítulo/sección: Dejamos que geometry actúe.
  }{
    % Es una figura: Neutralizamos geometry para que no las recorte mal.
    \geometry{pass}
  }
\fi
\makeatother
\usepackage{parskip} % Espaciado entre párrafos sin sangría
\usepackage{float}   % Para usar [H] en figura

% Ajustes de flotantes para evitar espacios verticales exagerados
\renewcommand{\topfraction}{0.95}      % Permite que las figuras ocupen hasta el 95% superior
\renewcommand{\bottomfraction}{0.95}   % Permite que las figuras ocupen hasta el 95% inferior
\renewcommand{\textfraction}{0.05}     % Permite hojas con tan solo un 5% de texto
\renewcommand{\floatpagefraction}{0.8} % Requiere que una página de solo figuras esté 80% llena
\setcounter{topnumber}{4}
\setcounter{bottomnumber}{4}
\setcounter{totalnumber}{10}

% --- 4. GRÁFICOS Y TABLAS ---
\usepackage{graphicx}
\usepackage{booktabs} % Tablas profesionales
\usepackage{caption}  % Control de captions y \captionof
\usepackage{tikz}
\usetikzlibrary{arrows.meta, calc, angles, quotes, babel, backgrounds}
\usepackage{pgfplots}
\usepgfplotslibrary{groupplots}
\usepgfplotslibrary{external}
\usepgfplotslibrary{patchplots}
\pgfplotsset{compat=1.18}
\usepackage[american, straightvoltages]{circuitikz}

% --- 5. COLORES Y CAJAS ---

\usepackage[table]{xcolor}
\usepackage{tcolorbox}

% Definición de colores corporativos/estilo
\definecolor{utnblue}{RGB}{0, 51, 102}
\definecolor{codegreen}{rgb}{0,0.6,0}
\definecolor{codegray}{rgb}{0.5,0.5,0.5}
\definecolor{codepurple}{rgb}{0.58,0,0.82}
\definecolor{backcolour}{rgb}{0.96,0.96,0.96}

% --- 6. ENTORNOS PERSONALIZADOS (CAJAS) ---

% Caja "Resultado" (Azul Corporativo) — Definiciones, fórmulas clave, resultados importantes.
% Uso: \begin{resultado}[Fórmula de Euler] ... \end{resultado}
\newtcolorbox{resultado}[1][]{
    colback=utnblue!5!white,
    colframe=utnblue,
    title=\textbf{#1},
    fonttitle=\bfseries,
    sharp corners,
    boxrule=0.5mm
}

% Caja "Nota" (Gris) — Advertencias, aplicaciones, contexto secundario.
% Uso: \begin{nota}[Cuidado con la calculadora] ... \end{nota}
% Uso: \begin{nota} ... \end{nota}  → título por defecto "Nota"
\newtcolorbox{nota}[1][Nota]{
    colback=codegray!5!white,
    colframe=codegray,
    title=\textbf{#1},
    fonttitle=\bfseries,
    sharp corners,
    boxrule=0.5mm
}

% Caja inline para resaltar un valor y enlazarlo con su otra aparición en una tabla
% o en una cuenta: mismo color, mismo valor.
% Uso: \marcavalor{$4$}  o  \marcavalor[codegreen]{$4$}
\newtcbox{\marcavalor}[1][utnblue]{
    on line,
    boxsep=1pt, left=2pt, right=2pt, top=1pt, bottom=1pt,
    colback=#1!10!white,
    colframe=#1,
    boxrule=0.4pt,
    arc=2pt
}

% --- 7. CÓDIGO FUENTE (LISTINGS) ---
\usepackage{listings}

\lstdefinestyle{miestilo}{
    backgroundcolor=\color{backcolour},   
    commentstyle=\color{codegreen},
    keywordstyle=\color{magenta},
    numberstyle=\tiny\color{codegray},
    stringstyle=\color{codepurple},
    basicstyle=\ttfamily\footnotesize,
    breakatwhitespace=false,         
    breaklines=true,                 
    captionpos=b,                    
    keepspaces=true,                 
    numbers=left,                    
    numbersep=5pt,                  
    showspaces=false,                
    showstringspaces=false,
    showtabs=false,                  
    tabsize=2,
    frame=single,                    
    rulecolor=\color{codegray}
}

\lstset{style=miestilo}
% Soporte para tildes y ñ en bloques de código
\lstset{literate={ñ}{{\~n}}1 {á}{{\'a}}1 {é}{{\'e}}1 {í}{{\'i}}1 {ó}{{\'o}}1 {ú}{{\'u}}1}

% Operadores matemáticos personalizados

% Vector en sentido discreto (lista finita de coordenadas): negrita + flecha.
% Uso: \vect{v}, \vect{e}_k, \vect{0}.
\newcommand{\vect}[1]{\vec{\mathbf{#1}}}

% Fecha de versión y comando para capítulos standalone
\providecommand{\verdate}{\the\year.\ifnum\month<10 0\fi\the\month.\ifnum\day<10 0\fi\the\day}
\newcommand{\chapterversion}{%
    \vspace{-1.5cm}\begin{flushright}\small\textit{\color{codegray}Versión~\verdate}\end{flushright}\vspace{0.5cm}%
}

% --- 8. HIPERVÍNCULOS (DEBE SER EL ÚLTIMO PAQUETE) ---
\usepackage[colorlinks=true, linkcolor=utnblue, urlcolor=utnblue, citecolor=utnblue]{hyperref}

% --- 9. REFERENCIAS CRUZADAS ENTRE CAPÍTULOS STANDALONE ---
% Cuando se compila un capítulo standalone, xr-hyper lee los .aux de los
% otros capítulos (si existen) para resolver \ref{} a labels externos.
% En la compilación del libro completo este bloque no se activa.
\ifstandalone
  \usepackage{xr-hyper}
  \makeatletter
  \newcommand{\xrchap}[1]{%
    \edef\xrc@self{\detokenize{Capitulo#1}}%
    \edef\xrc@job{\jobname}%
    \ifx\xrc@self\xrc@job\else
      \IfFileExists{../#1capitulo/Capitulo#1.aux}%
        {\externaldocument{../#1capitulo/Capitulo#1}}{}%
    \fi
  }
  \makeatother
  \xrchap{01}\xrchap{02}\xrchap{03}\xrchap{04}
  \xrchap{05}\xrchap{06}\xrchap{07}\xrchap{08}
  \xrchap{09}\xrchap{10}\xrchap{11}\xrchap{12}
  \xrchap{13}
\fi

\begin{document}
    \begin{tikzpicture}[>=Latex, font=\small, auto]
        % Estilos
        \tikzset{
            grid_x/.style={utnblue, thin},
            grid_y/.style={red!80!black, thin},
            axis_style/.style={->, black, thick},
            dot/.style={circle, fill=black, inner sep=1.5pt}
        }

        % --- PANEL A: PLANO Z (Rectangular) ---
        % Escala aumentada para equiparar áreas visuales. 
        % Dominio Z: x en [0, 1], y en [0, pi/2 approx 1.57]
        % Área matemática: 1 * 1.57 = 1.57
        % Vamos a escalar por 2.5
        \begin{scope}[scale=2.5]
            \node[anchor=south] at (0.5, 1.7) {\textbf{Plano Z}};
            
            % Ejes
            \draw[axis_style] (-0.2,0) -- (1.5,0) node[right] {Re ($x$)};
            \draw[axis_style] (0,-0.2) -- (0,1.8) node[above] {Im ($y$)};

            % Rejilla Vertical (x = c) -> Se mapearán a círculos
            \foreach \x in {0.2, 0.4, 0.6, 0.8, 1.0} {
                \draw[grid_x] (\x, 0) -- (\x, 1.6);
            }
            \node[utnblue, scale=0.8, anchor=north] at (1.0, 0) {$x=1$};

            % Rejilla Horizontal (y = c) -> Se mapearán a rayos
            % Usamos pasos de 0.3 radianes para que se vea claro
            \foreach \y in {0.3, 0.6, 0.9, 1.2, 1.5} {
                \draw[grid_y] (0, \y) -- (1.2, \y);
            }
            \node[red!80!black, scale=0.8, anchor=east] at (0, 1.5) {$y=1.5$};
            
            % Punto de interés: x=0.6, y=0.9
            \coordinate (z_pt) at (0.6, 0.9);
            \node[dot] at (z_pt) {};
            
            % Ángulo recto
            \draw[thin] (0.6, 1.0) -- (0.7, 1.0) -- (0.7, 0.9);
            
            % Etiqueta
            \node[anchor=south west, scale=0.7] at (0.7, 0.9) {$90^\circ$};
        \end{scope}

        % Flecha de transformación
        \draw[->, ultra thick, gray!50, bend left=15] (3, 2) to node[midway, above, black] {$w = e^z$} (5, 2);

        % --- PANEL B: PLANO W (Polar) ---
        % Dominio W: r en [1, e=2.72], theta en [0, 1.57]
        % Área matemática aprox: pi/4 * (2.72^2 - 1^2) approx 5
        % Escala para igualar: 2.5 / sqrt(5/1.57) approx 2.5 / 1.78 = 1.4
        \begin{scope}[shift={(6,0)}, scale=1.4]
            \node[anchor=south] at (1.5, 3.0) {\textbf{Plano W}};

            % Ejes
            \draw[axis_style] (-0.5,0) -- (3.5,0) node[right] {$u$};
            \draw[axis_style] (0,-0.5) -- (0,3.5) node[above] {$v$};

            % Círculos (Mapeo de x = c => |w| = e^x)
            \foreach \x in {0.2, 0.4, 0.6, 0.8, 1.0} {
                \pgfmathsetmacro{\rad}{exp(\x)}
                \draw[grid_x] (\rad, 0) arc (0:95:\rad);
            }
            % Etiqueta para r = e^1
            \pgfmathsetmacro{\radmax}{exp(1.0)}
            \node[utnblue, scale=0.8] at (2.0, -0.3) {$|w|=e$};

            % Rayos (Mapeo de y = c => arg(w) = y rad)
            \foreach \y in {0.3, 0.6, 0.9, 1.2, 1.5} {
                \pgfmathsetmacro{\deg}{\y * 180 / 3.14159}
                % Dibujamos rayos hasta un radio un poco mayor que e^1
                \draw[grid_y] (0,0) -- (\deg:3.0);
            }
            % Etiqueta para el último rayo
            \pgfmathsetmacro{\degmax}{1.5 * 180 / 3.14159}
            \node[red!80!black, scale=0.8] at (\degmax:3.2) {$\phi=1.5$};
            
            % Punto de interés mapeado: w = e^(0.6) /_ 0.9 rad
            \pgfmathsetmacro{\rw}{exp(0.6)}
            \pgfmathsetmacro{\degw}{0.9 * 180 / 3.14159}
            \coordinate (w_pt) at (\degw:\rw);
            
            \node[dot] at (w_pt) {};
            
            % Ángulo recto sobre la intersección curva
            % Tangente del circulo: dirección degw + 90
            % Tangente del rayo: dirección degw
            \begin{scope}[shift={(w_pt)}, rotate=\degw]
                \draw[thin] (0.2, 0) -- (0.2, 0.2) -- (0, 0.2);
                \node[scale=0.7, anchor=south west, rotate=-\degw] at (0.2, 0.2) {$90^\circ$};
            \end{scope}

        \end{scope}

    \end{tikzpicture}
\end{document}
